\lesson{09/10}
[...] qualcosa su T fissata, $p = p(T)\pm\varepsilon$ usando $G$ si vede che c'è
un $\Delta V$

How can the shape of the transition curve $p(T)$ be determined? With the \textit{
Clausius-Clapeyron equation}.

Recall that
$$\mu(T,p) = \frac{G}{N}\quad \dd\mu = - s \dd T + v\dd p$$
and this has to be satisfied at each point on the curve: consider the
infinitesimal change in $\mu$ along the curve
\begin{gather*}
    \mu_\alpha \rightarrow \mu_\alpha+\dd\mu_\alpha\\
    \mu_\beta \rightarrow \mu_\beta+\dd\mu_\beta
\end{gather*}
since $\mu_\alpha = \mu_\beta$ at all points,
$$\dd\mu_\alpha = \dd\mu_\beta$$
substituting yields
$$-s_\alpha\dd T + v_\alpha\dd p = -s_\beta\dd T + v_\beta \dd p$$
which implies
$$\dv[]{p}{T} = \frac{\Delta s}{\Delta v} = \frac{\Delta S}{\Delta V}$$
in terms of the latent heat $L = T\Delta S$ it can be written as
$$\dv[]{p}{T} = \frac{L}{T\Delta V} = \frac{1}{T}\frac{l}{\Delta v}=
\frac{1}{T}\frac{l_\text{mol}}{\Delta v_\text{mol}} = \frac{1}{T}\frac{l_\text{sp}}
{\Delta v_\text{sp}}$$
where $l = \frac{L}{N}$ is the latent heat per particle, $l_\text{mol}$ is
the molar latent heat and $l_\text{sp}$ is the specific latent heat (J/kg).

\textbf{Example} water to ice transition\\
In standard conditions, $p_\text{atm} = 101325$ Pa, $T_0 = 273.15$ K, the transition
is possible and, contrary to many substances, water is denser than ice: the variation
of specific volume for the transition is
$$\Delta v_\text{sp} = v^\text{ice}_\text{sp} - v^\text{water}_\text{sp} = 9\cdot
10^{-5}\text{ m}^3/\text{kg}$$
and the specific latent heat is
$$l_\text{sp} = 2.3\cdot10^{5}\text{ J/kg}$$
The question is, how much pressure is needed to melt the ice at temperature
$T_0-1$ K?

In this case, $1\text{ K} =\Delta T << T_0$ and the C-C equation becomes
$$\dv[]{p}{T} = \frac{1}{T}\frac{l_\text{sp}}{\Delta v_\text{sp}}\rightarrow
\frac{\Delta p}{\Delta T} = \frac{1}{T}\frac{l_\text{sp}}{v_\text{sp}}\rightarrow
\Delta p = \frac{\Delta T}{T_0} \frac{l_\text{sp}}{\Delta v_\text{sp}}\approx 10^7
\text{ Pa}
$$

\textbf{Example} vapour at saturation\\
When the liquid and gas phases coexist, the vapour is at saturation; this is a first
order transition in which $l>0$ and $v_\text{gas}>>v_\text{liquid}$, so $\Delta v 
\approx v_\text{gas}$. The C-C equations takes form as
$$\dv[]{p}{T} = \frac{l}{Tv_\text{gas}}$$
assuming a perfect gas that satisfies the ideal state equation
$$pV = NkT\rightarrow pv_\text{gas} = kT$$
the equation becomes
$$\dv[]{p}{T} = \frac{l}{T}\frac{p}{kT}\rightarrow \frac{\dd p}{p} = \frac{l}{k}
\frac{\dd T}{T^2}\rightarrow\ln\left(\frac{p}{p_0}\right) = -\frac{l}{k}\left(
\frac{1}{T}-\frac{1}{T_0}\right)$$
so the form of the transition curve is an exponential
$$p(T) = p_0\exp\left[-\frac{l}{k}\left(\frac{1}{T}-\frac{1}{T_0}\right)\right]$$

What about the $pV$ plane? Since the control variables can always be chosen, it should
not be fundamentally different: consider a coexistence curve with a point
$(p(T_0),T_0)$ and $\dv[]{p}{T}(T_0)>0$ and a transition at constant $T$ and
decreasing $p$; the transition follows an isothermal curve in the $pV$ plane, then
at $p(T_0)$ it becomes "isobara" while the volume changes, this is the coexistence
curve, and finally it becomes isothermal again.\\
Consider the coexistence curve: during the "isobara"
$$V = N_\alpha v_\alpha \underset{\text{"isobara"}}\longrightarrow V = N_\alpha
v_\alpha + N_\beta v_\beta\rightarrow v = \nu_\alpha v_\alpha + \nu_\beta v_\beta$$
where $\nu_i = N_i/N$; from this follows
$$v-v_\alpha = \nu_\beta(v_\beta-v_\alpha)$$
and
$$v_\beta - v = \nu_\alpha(v_\beta-v_\alpha)$$
their quotient yields
$$\frac{v-v_\alpha}{v_\beta-v} = \frac{\nu_\beta}{\nu_\alpha}$$

Considering all the different values of $T$ shows that in the $pV$ plane
the coexistence happens in a 2D region instead of a curve.

$v_\alpha$ and $v_\beta$ are implicitly determined by
$$\begin{dcases*}
  p_\alpha = p_\beta\\
  \mu_\alpha = \mu_\beta
  \end{dcases*}$$
In terms of free energy per particle $f = F/N$, $p = -\pdv[]{f}{v}$:
since $p$ is constant in the phase transition, $f(v)$ is linearly decreasing with
slope
$$-p = \frac{f_\alpha - f_\beta}{v_\beta-v_\alpha}$$ 
so
$$f_\beta-f_\alpha = -p (v_\beta-v_\alpha) \rightarrow (f+pv)_\alpha = (f+pv)_\beta$$
recalling that $F$ is related with a Legendre transform to the Gibbs free energy $G$,
$$f+pv = G/N = \mu$$
it implies
$$\mu_\alpha = \mu_\beta.$$
It can be shown that
$$f(v) = \nu_\alpha f_\alpha + \nu_\beta f_\beta$$
\subsection{Continuous transitions}
In this case the first derivatives of $\mu$ are continuous and that means
$$\Delta S = 0\quad\Delta V = 0$$
and the discontinuity is at higher orders. Typically the coexistence in the
plane $p,T$ is represented by single points, as opposed to curves for first
order transitions: these are called \textit{critical points} and usually they
are endpoints of first order transitions.

\section{Maxwell construction}
Suppose to know an approximated empirical equation of state for some system
$$p=p(V,T)$$
that in some regions \textit{does not} obey the stability conditions, i.e.
in some regions, at constant temperature,
$$\eval{\pdv[]{p}{V}}_{T,N}>0$$
this is not possible and means that the equation of state is wrong, but suppose it
works well in other regions: how can this be corrected?\\
There is a method to connect "good" regions with one another, due to Maxwell, that
consists in supposing that the "bad" region represents a phase transition: for
example connecting two points in the $pV$ plane with a straight line at $p=\bar p$;
this connection must be between two points $V_1,V_2$ so that the free energy at
those points is the same as the one determined by the equation of state.
$$\dd F = -S\dd T -p\dd V +\mu \dd N$$
fixing $T,N$ this becomes
$$F(V_2,T,N) - F(V_1,T,N) = -\int_{V_1}^{V_2}\dd V p(V,T,N)$$
On the contrary, for the "isobara" the difference in free energy is
$$F(V_2,T,N)-F(V_1,T,N)=\bar p (V_2-V_1)$$
so the condition is
$$\bar p (V_2-V_1) = \int_{V_1}^{V_2}\dd V p(V,T,N)$$  
\textbf{Example: Van Deer Waals}\\
The VdW equation considers intermolecular interactions with a potential "wall" at
distance $d$ and for $r>d$ a $\sim -\frac{1}{r}$ potential: this is a good
approximation of the real potential that governs intermolecular interactions;
this also determines an effective volume
$$V_\text{eff} = V-b$$
as the single molecule's "volume" is $\sim d^3$ and this volume is unavailable for
the rest of the system. The well after distance $d$ forms bound states, which 
lowers the kinetic pressure:
$$p = p_\text{kin} - \frac{a}{V^2}$$
this defect in pressure is proportional to $V^{-2}$ (why?).

The last hypothesis is that the kinetic pressure and the effective volume follow
the ideal gas law
$$p_\text{kin}V_\text{eff} = nRT$$
substituting yields the VdW equation of state
$$\left(p+\frac{a}{V^2}\right)(V-b) = nRT$$
since $p(V)$ is approximately dependent on $V^{-3}$, some isothermals are
non-physical, since they present horizontal flex points.
